\documentstyle[12pt,a4]{article}
\newtheorem{lema}{\bf Lemma}[section]
\newtheorem{corolario}[lema]{\bf Corolary }
\newtheorem{teorema}[lema]{\bf Theorem}
\newtheorem{ejemplo}[lema]{\bf Example}
\newtheorem{nota}[lema]{\bf Remark}
\def\demostracion{\noindent{\bf Proof: }}
%\def\gg{{\frak g}}
\def\gg{{\got g}}
\def\CC{{\Bbb C}}
\def\KK{{\Bbb K}}
\def\NN{{\Bbb N}}
\def\RR{{\Bbb R}}
\def\ZZ{{\Bbb Z}}
\title{The use of $Mathematica$ for the
classification of some nilpotent Lie algebras}
\author {L.M. Camacho\footnotemark[1]
\and J.R. G\'{o}mez\thanks{Dpto. Matem\'{a}tica Aplicada I. Univ.
Sevilla. Avda. Reina Mercedes S.N. 41012 Sevilla (Spain).
E-mail: jrgomez@cica.es, lcamacho@euler.fie.us.es}
\and R.M. Navarro\thanks{Dpto. de Matem\'{a}ticas.
Univ. de Extremadura. Avda. de la Universidad S.N.
10071 C\'{a}ceres (Spain). E-mail: rnavarro@unex.es}}
\date{}
\begin{document}
\maketitle
\begin{abstract}
The classification of nilpotent Lie algebras is
nowadays an open problem \cite{biblia}. In the last years have
been studied the nilpotent Lie algebras leading to
their classification up to dimension $7$, as a
particular type of the above, the classification
of the filiform Lie algebras is giving up to
dimension $11$ and the classification of the
2-abelian filiform Lie algebras for the case of
generic dimension. Cabezas and G\'{o}mez generalize
the notion of filiform algebra to $p$-filiform
algebra \cite{articulo kluwer}, \cite{Comm(n-4)},
\cite{extracta} (whose Goze's invariant is
$(dim(\gg)-p,1,1,\dots,1)$), also give the
classification for the integer values of $p$
between $dim(\gg)-4$ and $dim(\gg)-2$, for
$p=dim(\gg)-5$ we have classified the
$(n-5)$-filiform Lie algebras with maximal derived
algebra, where $n=dim(\gg)$, \cite{jljr}.
\vglue2.5mm
The determination of Lie algebras is a difficult
work, therefore we have developed some new methods
which allows the classification of nilpotent Lie
algebras.
%Many papers are based on direct
%computations (by hand) and the complexity of
%those computations leads frequently to errors.
\vglue2.5mm
In this paper, we show a program by using the
$Mathematica\ 3.0$ package which allows the
classification of the sets of laws of
$(n-6)$-filiform Lie algebras of dimension $8$.
\end{abstract}
\begin{section}{Description of the program}
Classifications problems up to isomorphisms can
be reduced to determinate the structure constants
of non-isomorphic algebras. As the structure
constants depend on the basis, it is convenient
to find a basis for which one has a maximum
number of structure constants equal zero. We
consider adapted basis.

Let be a family of $(n-6)$-filiform Lie algebras
of dimension $8$. For the classification of this
family, it is necessary to find classic
invariants. This is not always possible. Then, it
is necessary to make the generic change of basis.
This program allows to generate the new vectors
and to compute new structure constants. Later we
have to do a discussion over the output of program
to can classify.
\vglue2mm
The way to find non-isomorphic algebras of
dimension $8$ can be summarized in the following
steps
%\vglue1mm

Step 1: We introduce the family, the conditions of
bilinear and anti-symmetry and the condition of
$(n-6)$-filiform Lie Algebras.
\vglue 3mm
Step 2: We generate particular new adapted basis,
where the structure constants can be zero.
\vglue3mm
Step 3: We compute the conditions to
maintain some null bracket products. We select
only most simple null bracket products.
\vglue3mm
Step 4: We obtain new structure constants of the
family in the new adapted basis. In this
way, we compute some bracket products.
\vglue3mm
Step 5: We check that the remain of null bracket
products. In the new adapted basis these products
have to be null.
\vglue3mm
Finally, we discuss if the nullities of new structure
constants are invariants or zero. This discussion allows to
classify, since it separates in non-isomorphic
algebras.

\end{section}

\begin{thebibliography}{1}
%\bibitem{tesis Jesus}J.M. Cabezas, {\it Una generalizaci\'{o}n de las
%\'{a}lgebras de Lie filiformes,} Phd thesis, Universidad de Sevilla, 1996.

\bibitem{jljr} J.M. Cabezas, L.M. Camacho, J.R. G\'{o}mez, R.M. Navarro,
{\it A class of nilpotent Lie algebras,} Submitted
Communications in Algebra.

\bibitem{extracta}  J.M. Cabezas, J.R. G\'{o}mez,
{\it  Las \'{a}lgebras de Lie $(n-3)$-filiformes como
extensiones por derivaciones,} Extracta Mathematicae,
13:3, 383-391, 1999.


\bibitem{Comm(n-4)}  J.M. Cabezas, J.R. G\'{o}mez, {\it  (n-4)-filiform
Lie algebras,} Communications in Algebra,
27:10, 1999.

\bibitem{articulo kluwer}
J.M. Cabezas, J.R. G\'{o}mez, A. Jim\'{e}nez-Merch\'{a}n, {\it Family of p-filiform
Lie algebras,} In Algebra and Operator Theory, 1997. Ed. Y. Khakimdjanov,
M. Goze, Sh. Ayupov, 93-102, Kluwer Academic Publishers, 1998.


\bibitem{biblia} M. Goze, Y. Khakimdjanov, {\it Nilpotent Lie
Algebras}, Kluwer Academic Publishers, 1996.

\end{thebibliography}
\end{document}

